\begin{table}[H]
\centering
\caption{Remuneración de ocupados de 25 a 29 años con pregrado universitario, 2010 y 2025}
\label{tab:remuneracion_pregrado_universitario_25_29_larga}
\footnotesize
\begin{tabular}{@{}p{7.4cm}rr@{}}
\toprule
Indicador & 2010 & 2025 \\
\midrule
\multicolumn{3}{@{}l}{\textbf{Panel A. Remuneración mensual por trabajador}} \\
Ocupados (miles) & 228,1 & 493,8 \\
Promedio (millones de pesos de 2025) & 3,3 & 3,0 \\
Mediana (millones de pesos de 2025) & 2,5 & 2,5 \\
Desviación estándar (millones de pesos de 2025) & 3,2 & 2,6 \\
Razón P90/P10 & 5,2 & 3,5 \\
Entre 0,9 y 1,1 SMLMV & 5,7\% & 12,6\% \\
\addlinespace[0.6em]
\multicolumn{3}{@{}l}{\textbf{Panel B. Remuneración por hora trabajada}} \\
Promedio (miles de pesos de 2025) & 17,6 & 16,3 \\
Mediana (miles de pesos de 2025) & 13,8 & 13,3 \\
Desviación estándar (miles de pesos de 2025) & 16,6 & 14,4 \\
Razón P90/P10 & 5,0 & 4,0 \\
\bottomrule
\end{tabular}
\caption*{\footnotesize Nota: la muestra se restringe a ocupados de 25 a 29 años cuyo logro educativo armonizado es pregrado universitario. La GEIH no identifica el año de graduación. Por lo tanto, este grupo aproxima a los recién egresados, pero no los observa directamente. El SMLMV se expresa en pesos constantes de 2025 y no incluye auxilio de transporte. Fuente: cálculos propios con GEIH del DANE.}
\end{table}
